\begin{finding}[\fid{F27}: what changing parity or the effective cone accomplishes]\label{f:derived}
Let $R=\Z[v,v^{-1}]$, $\delta=v+v^{-1}$,
$H=R[c]/(c^2-\delta c)$, and let $P$ be the actual rank-one element
of \fid{F09}. In a graded additive monoidal categorification whose
split Grothendieck ring has the named effective basis
$1,c_1,c_2,c_1c_2$, the following hold.
\begin{enumerate}
\item $P$ has an explicit bounded two-term lift, but no bounded complex
with Euler class $P^k$ ($k\ge1$) has all its nonzero terms in a single
homological parity. An overall sign or internal grading shift does not
alter this obstruction.
\item The explicit zero-differential lift has tensor-power homotopy
amplitude $k$. Separately, if a nonzero bounded complex in an abelian tensor category
with faithful exact vector-space fiber functor has extreme nonzero
cohomology degrees $p,r$, its $k$th tensor power has amplitude $k(r-p)$.
Mixed cohomological parity persists in every tensor power.
\item A different positive cone does preserve the actual rank-one
character actions. The faithful map
\[
 H\longrightarrow R^2,\qquad a+bc\longmapsto(a+b\delta,a)
\]
becomes an isomorphism after inverting $\delta$. Its tensor square sends
\[
 P\longmapsto(\delta^2,0,0,\delta^2)
\]
on the four labels $++,+-,-+,--$. Thus the missing sign-tensor-sign
action in \fid{F16} is restored. In the even form $u=v^2$, the two
nonzero values are actual positive Tate classes
$u^{-1}+2+u$.
\end{enumerate}
These statements do not identify the new character cone with VIII's
canonical basis or construct the negative path-output spectrum of
\fid{F24}.
\end{finding}
\begin{proof}
The already proved quadratic relation gives
\[
 P^k=\delta^{2k-2}P
 =\delta^{2k}\,1-\delta^{2k-1}(c_1+c_2)
                  +2\delta^{2k-2}c_1c_2.
\]
Its four named coordinates have both signs for every $k$.
Let $B_0,B_1,B_2,B_3$ represent $1,c_1,c_2,c_1c_2$, with $B_0$ the
tensor unit. Internal shift $(1)$ represents $v$ and is distinct
from homological shift $[1]$. Set
\[
 X_+=B_3^{\oplus2}\oplus B_0(2)\oplus B_0^{\oplus2}\oplus B_0(-2),
\]
\[
 X_-=B_1(1)\oplus B_1(-1)\oplus B_2(1)\oplus B_2(-1).
\]
The complex $K^0=X_+$, $K^1=X_-$ with zero differential has class
$P$. Any complex with terms only in even degrees has effective Euler
class, and one with terms only in odd degrees has the negative of an
effective class. Neither can equal the displayed mixed-sign class of
$P^k$. Homotopy equivalence preserves Euler class, so it cannot repair
this; overall signs and grading monomials preserve the mixed pattern.

The zero-differential $K^{\otimes k}$ has nonzero terms in all degrees
$0,\ldots,k$. Use unit summands from $X_+$ and a $B_1$ summand from
each chosen $X_-$; $c_1^j=\delta^{j-1}c_1\ne0$ proves nonvanishing
when $j\ge1$. If a zero-differential complex is homotopy equivalent
to $D$, the composite maps back to it equal the identity degreewise,
because its homotopy correction $dh+hd$ is zero. Hence its degree-$i$
term is a retract of $D^i$; no nonzero extreme can disappear. This
proves the amplitude claim for this particular lift, not for every
complex with its Euler class.

For the separate cohomological assertion, apply the faithful exact
tensor fiber functor. Over a field the K\"unneth decomposition of
$H^j(K^{\otimes k})$ is the direct sum of
$H^{i_1}(K)\otimes\cdots\otimes H^{i_k}(K)$ with sum of degrees $j$.
The extreme summands are the nonzero powers of $H^p$ and $H^r$.
Opposite-parity nonzero groups in degrees $i,j$ give nonzero summands
in degrees $ki$ and $(k-1)i+j$, of opposite parity. Exactness and
faithfulness detect these assertions in the original category.
Shifts, Tate twists and finite base-field extensions do not reduce
the amplitude. This argument does not assume an arbitrary convolution
or pushforward is exact for a chosen heart.

For the character model, the map evaluates $c$ at $\delta$ and zero.
Its kernel is zero since $R$ is a domain and $\delta\ne0$; its image
is the pairs whose difference is in $\delta R$. Localizing makes it
surjective. The four-coordinate tensor map is injective by successive
evaluation of $1,c_1,c_2,c_1c_2$, and direct evaluation of $P$ gives
the stated tuple. At $v=1$, $c_i=1+s_i$ and $P=2(1+s_1s_2)$,
so both equal-sign characters have eigenvalue four and both others
have eigenvalue zero. Every module over the localized algebra splits
by its four actual orthogonal character idempotents. It becomes a
quasi-coherent sheaf on four copies of the localized base, a vector
bundle when its corresponding modules are finite projective.

For an integer-Tate formulation set $b=vc$ and use
\[
 H_{\rm even}=\Z[u,u^{-1}][b]/(b^2-(1+u)b),\qquad
 P=u^{-1}\{b_1b_2+(b_1-1-u)(b_2-1-u)\}.
\]
Scalar extension recovers the original algebra. On each rational
character point, represent $b_i$ by $\mathbf1\oplus\mathbf1(-1)$
if its sign is $+$, and by zero otherwise. The element $P$ is then
represented on each diagonal point by
\[
 T=\mathbf1(1)\oplus\mathbf1^{\oplus2}\oplus\mathbf1(-1),
 \qquad \Tr(F^n\mid T)=q^{-n}+2+q^n,
\]
and by zero off the diagonal. Tensoring with this object on the
two supported points gives the corresponding positive operation on
the character Grothendieck group. It has weights $-2,0,0,2$;
it is not a single pure weight-zero representation.
\end{proof}

There is an actual Hecke category satisfying the first hypothesis:
the product of the two rank-one Soergel categories. For one factor
take $A=\Q[x]$, $\deg x=2$, $s(x)=-x$, and
$B=A\otimes_{A^s}A(1)$. The rank-two free decomposition over
$A^s=\Q[x^2]$ gives $B\otimes_A B=B(1)\oplus B(-1)$.
The split Grothendieck/KL identification is the known result
\cite[\S3.5, Theorem~3.5, pp.~1107--1108; Theorem~1.1, p.~1093]{EW}.
The bounded category used above is its homotopy category; no
unspecified perverse $t$-exactness was invoked.

The character map is also a finite normalization of affine schemes:
$R^2$ is the integral closure of $H$ in its total fraction ring,
and $R^4$ is that of the tensor-square algebra. Indeed the separate
branch idempotents are integral; conversely an element integral over
these finite $R$-algebras is integral over the normal ring $R$ and
has each coordinate in $R$. The map is an isomorphism on
$\delta\ne0$, but it is not an integral identification of the named
KL effective basis with the full character basis.

The finite-spectrum lemma of \fid{F24} shows why tensoring the
derived input does not itself manufacture the required opposite sign:
every nonempty visible real Euler spectrum has a positive top-radius
coefficient in each even tensor power. Input shifts contribute sign
$(-1)^{kr}=1$ there. An extra odd shift of the output reverses a sign
formally, but its identification with path-hypersurface cohomology
still requires the trace identity. Full scalar localization also
cannot be silently realized by finite Weil traces:
$(1+q^n)^{-1}$ is excluded by \fid{F26}. A particular localized
idempotent may have its denominator cancel in character coordinates,
as $b/(1+u)\mapsto(1,0)$ does. The obstruction concerns the specified
non-Laurent all-extension trace, not every object in a localized
category.


